\begin{table}[htbp]
\centering
\caption{Long-Short Portfolio Composition (1963--2024)}
\label{tab:longshort_composition}

\begin{adjustbox}{width=0.9\linewidth}
\begin{tabular}{l l*{5}{r}}
\toprule
\textbf{Category} & \textbf{Signal} & \textbf{Avg P1} & \textbf{Avg P10} & \textbf{Total} & \textbf{P1\%} & \textbf{P10\%} \\
\midrule
\multicolumn{7}{l}{\textit{Panel A: Profitability}} \\
 & Gross Profitability & 47.0 & 52.7 & 99.7 & 47.2 & 52.8 \\
 & Operating Profitability & 53.0 & 58.0 & 111.0 & 47.8 & 52.2 \\
 & Return on Assets & 52.3 & 66.2 & 118.5 & 44.2 & 55.8 \\
 & $\Delta$ Asset Turnover & 49.7 & 50.2 & 99.9 & 49.8 & 50.2 \\
 & $\Delta$ Profit Margin & 52.6 & 55.2 & 107.7 & 48.8 & 51.2 \\
\addlinespace
\multicolumn{7}{l}{\textit{Panel B: Investment}} \\
 & Return on Invested Capital & 49.2 & 64.3 & 113.5 & 43.3 & 56.7 \\
 & Asset Growth & 41.0 & 62.6 & 103.6 & 39.6 & 60.4 \\
 & Investment (CapEx) & 48.7 & 54.9 & 103.5 & 47.0 & 53.0 \\
 & CapEx + Inventory Growth & 37.7 & 51.7 & 89.4 & 42.1 & 57.9 \\
 & Net Share Issuance & 53.1 & 63.0 & 116.1 & 45.7 & 54.3 \\
 & CapEx Growth & 48.8 & 58.0 & 106.8 & 45.7 & 54.3 \\
\addlinespace
\multicolumn{7}{l}{\textit{Panel C: Accruals \& Quality}} \\
 & Sales Growth & 41.3 & 65.8 & 107.1 & 38.6 & 61.4 \\
 & Working Capital Accruals & 44.8 & 52.8 & 97.6 & 45.9 & 54.1 \\
 & Sales-Receivables Gap & 41.0 & 55.3 & 96.3 & 42.6 & 57.4 \\
 & Receivables Turnover & 55.2 & 46.6 & 101.8 & 54.2 & 45.8 \\
\addlinespace
\multicolumn{7}{l}{\textit{Panel D: Liquidity \& Solvency}} \\
 & Operating Leverage & 40.2 & 58.9 & 99.1 & 40.6 & 59.4 \\
 & Cash Cushion & 42.0 & 87.3 & 129.4 & 32.5 & 67.5 \\
 & Book Leverage & 82.9 & 39.8 & 122.7 & 67.5 & 32.5 \\
 & $\Delta$ Current Ratio & 56.3 & 57.7 & 114.0 & 49.4 & 50.6 \\
 & Debt Capacity & 74.6 & 44.0 & 118.5 & 62.9 & 37.1 \\
\addlinespace
\multicolumn{7}{l}{\textit{Panel E: Operating Efficiency}} \\
 & Cash Productivity & 75.7 & 39.4 & 115.1 & 65.8 & 34.2 \\
 & Depreciation/PPE & 43.0 & 68.6 & 111.5 & 38.5 & 61.5 \\
 & $\Delta$ Depreciation & 43.6 & 61.2 & 104.8 & 41.6 & 58.4 \\
 & Tax/Income Ratio & 55.8 & 48.0 & 103.8 & 53.8 & 46.2 \\
 & Asset Turnover & 55.0 & 48.9 & 103.9 & 52.9 & 47.1 \\
\midrule
\textbf{MEAN} & & \textbf{51.4} & \textbf{56.4} & \textbf{107.8} & \textbf{47.5} & \textbf{52.5} \\
\bottomrule
\end{tabular}
\end{adjustbox}

\medskip

\begin{minipage}{\textwidth}
\footnotesize
\textbf{Notes:} This table reports the average composition of long-short portfolios constructed from decile sorts on 25 fundamental signals. \textit{Avg P1} (P10) is the time-series average number of firms in the lowest (highest) decile portfolio. \textit{Total} is the sum of P1 and P10. \textit{P1\%} (P10\%) is the percentage of total firms allocated to each portfolio. Portfolios are formed annually in June using NYSE decile breakpoints applied to the filtered universe (NYSE P20 filter + price $\geq$\$5). Asymmetries in P1 vs. P10 reflect signal-specific data completeness constraints. For example, Book Leverage (67.5\% in P1 vs. 32.5\% in P10) indicates more firms with high leverage than low leverage, while Cash Cushion shows the opposite pattern (32.5\% vs. 67.5\%). Sample period: July 1963 to December 2024.
\end{minipage}

\end{table}
